%-----------PUBLICATIONS-----------------
\section{Research and Publications}
  \resumeSubHeadingListStart
    \resumeSubheading
      {Morphling: Emulator for Distributed Machine Learning at the Edge}{EdgeSys 2026}
      {Leyang Xue, \textbf{Yufeng Xia}, Eren Mendi, Ismaeel Bashir, Jiaxun Yang, Myungjin Lee, Mahesh Marina}{}
      \resumeItemListStart
        \resumeItem{Scalable Edge ML Emulator}
          {Developed Morphling, a measurement-driven emulator that evaluates distributed ML training on heterogeneous edge devices at scale, supporting up to 1800 emulated devices on a single GPU server with minimal latency and thermal error.}
      \resumeItemListEnd

    \resumeSubheading
      {Weaver: Foundation Model Training over AI-RAN Compute Infrastructure}{MobiCom 2026}
      {Leyang Xue, Tianxin Wang, Xin Zhe Khooi, Jiaxun Yang, Dheeraj Mahendiran, \textbf{Yufeng Xia}, Mun Choon Chan, Myungjin Lee, Mahesh Marina}{}
      \resumeItemListStart
        \resumeItem{AI-RAN Infrastructure Sharing}
          {Developed Weaver, a system enabling opportunistic foundation model training on GPU-accelerated RAN infrastructure, achieving 2.1--4.2$\times$ higher training throughput while harvesting up to 83\% of spare compute without impacting RAN performance.}
      \resumeItemListEnd

    \resumeSubheading
      {EmuRAN: Scalable and Cost-Effective Cloud-based RAN Emulation System}{MobiCom 2026}
      {Ujjwal Pawar, Andrew E. Ferguson, \textbf{Yufeng Xia}, Xenofon Foukas, Mahesh Marina, Bozidar Radunovic}{}
      \resumeItemListStart
        \resumeItem{Scalable RAN Emulation}
          {Developed EmuRAN, a cloud-based RAN emulation system capable of emulating 10,000 mobile devices and 10,000 base stations on 130 cloud instances, two orders of magnitude larger than existing solutions.}
      \resumeItemListEnd

    \resumeSubheading
      {Enhanced Pneumonia Detection in Chest X‑Rays Based on Integrated Denoising Autoencoders and Convolutional Neural Networks}{2024}
      {YUFENG XIA}{}
      \resumeItemListStart
        \resumeItem{Research Focus}
          {Conduct research on integrating machine learning with medical diagnostics to improve cancer classification accuracy using denoising autoencoders.}
      \resumeItemListEnd

    \resumeSubheading
      {CoSense: Auto‑Quantization Based on Sensor Specifications for Compiler Optimizations}{Under Review, 2025}
      {YUFENG XIA, PEI MU, ANTONIO BARBALACE}{}
      \resumeItemListStart
        \resumeItem{Framework Development}
          {Developed a novel LLVM‑based auto‑quantization framework leveraging sensor specifications to enhance performance and reduce code size for embedded systems.}
        \resumeItem{Quantization Techniques}
          {Implemented quantization techniques tailored to sensor data, achieving significant improvements in execution speed and binary optimization.}
        \resumeItem{Evaluation}
          {Evaluated CoSense on ARM Cortex‑M0 and Cortex‑A53 platforms, demonstrating measurable improvements in power, performance, and area (PPA).}
      \resumeItemListEnd

  \resumeSubHeadingListEnd